% Part: normal-modal-logic
% Chapter: filtrations
% Section: introduction

\documentclass[../../../include/open-logic-section]{subfiles}

\begin{document}

\olfileid[pt]{nml}{fil}{int}

\olsection{Introdução}

Uma questão importante acerca de uma lógica é sempre saber se ela é
decidível, isto é, se existe um procedimento efetivo que responda à
pergunta ``esta !!{formula} é válida?''. A lógica proposicional é
decidível: podemos testar efetivamente se !!a{formula} é uma tautologia
construindo uma tabela de verdade e, para uma !!{formula} dada, a
tabela de verdade é finita. Mas não podemos testar de maneira óbvia se
uma fórmula modal é verdadeira em todos os modelos, pois há
infinitos modelos. Podemos enumerar todos os modelos finitos relevantes
para uma fórmula dada, pois só são relevantes as atribuições de
subconjuntos de mundos às !!{propositional variable}s que de fato
ocorrem na !!{formula}. Se a relação de acessibilidade é fixa, as
diferentes atribuições possíveis $V(p)$ são simplesmente todos os
subconjuntos de~$W$ e, se $\card{W} = n$, há $2^n$ deles. Se nossa
!!{formula}~$!A$ contém $m$ !!{propositional
  variable}s, há então $2^{nm}$ modelos diferentes com~$n$ mundos.
Para cada um, podemos testar se~$!A$ é verdadeira em todos os mundos,
simplesmente calculando o valor de verdade de~$!A$ em cada um deles. É
claro que também devemos verificar todas as relações de acessibilidade
possíveis, mas há apenas um número finito de relações sobre~$n$ mundos
(mais precisamente, o número de subconjuntos de $W \times W$, isto é,
$2^{n^2}$.

Se não estamos interessados na lógica~\Log{K}, mas em uma lógica
definida por alguma classe de modelos (por exemplo, os modelos
reflexivos e transitivos), também devemos ser capazes de testar se a
relação de acessibilidade é do tipo correto. Podemos fazê-lo sempre que
os frames de nosso interesse são definíveis por !!{formula}s modais
(por exemplo, testando se \Ax{T} e \Ax{4} são válidos no frame). Assim,
a ideia seria percorrer todos os frames finitos, testar se cada um é um
frame da classe de nosso interesse, enumerar em seguida todos os
modelos possíveis sobre esse frame e testar se $!A$ é verdadeira em
cada um. Se não for, paramos: $!A$ não é válida na classe de modelos
de interesse.

Há um problema com essa ideia: não sabemos quando, se é que isso
ocorrerá, podemos parar de procurar. Se a !!{formula} tem um
contramodelo finito, nosso procedimento o encontrará. Mas, se ela não
tem contramodelo finito, não obteremos uma resposta. A !!{formula} pode
ser válida (sem qualquer contramodelo), ou pode ter apenas um
contramodelo infinito, que nunca examinaremos. Esse problema pode ser
superado se pudermos mostrar que toda !!{formula} que tem um
contramodelo tem um contramodelo finito. Nesse caso, dizemos que a
lógica tem a \emph{propriedade dos modelos finitos}.

Mas como mostrar que uma lógica tem a propriedade dos modelos finitos?
Uma maneira seria encontrar um modo de transformar um (contra)modelo
infinito de~$!A$ em um modelo finito. Se isso puder ser feito, sempre
que houver um modelo no qual $!A$ não é verdadeira, o modelo finito
resultante também tornará $!A$ não verdadeira. Esse modelo finito
aparecerá em nossa lista de todos os modelos finitos e, para toda
fórmula que não seja válida, acabaremos determinando que ela não é
válida. Nosso procedimento não terminará se a fórmula for válida. Se,
além disso, pudermos mostrar que existe algum tamanho máximo que o
modelo finito fornecido por nosso procedimento pode ter, e que esse
tamanho máximo depende apenas da !!{formula}~$!A$, teremos um tamanho
até o qual devemos testar modelos finitos em nossa busca por
contramodelos. Se até então não tivermos encontrado um contramodelo,
não haverá nenhum. Nosso procedimento decidirá então, de fato, a
pergunta ``$!A$ é válida?'' para qualquer !!{formula}~$!A$.

Uma estratégia que muitas vezes funciona para transformar estruturas
infinitas em estruturas finitas é ``identificar'' !!{element}s da
estrutura que se comportam da mesma maneira nos aspectos relevantes.
Se há infinitos mundos em~$\mModel{M}$ que se comportam da mesma
maneira nos aspectos relevantes, talvez possamos reunir \emph{todos}
os mundos em um número finito de ``classes'' (possivelmente infinitas)
de tais mundos. Em outras palavras, poderíamos tentar particionar o conjunto de
mundos da maneira correta, isto é, de modo que cada parte contenha
infinitos mundos, mas haja apenas um número finito de partes. Definimos
então um novo modelo~$\mModel{M^*}$ cujos mundos são as partes. Um
número finito de partes no modelo antigo nos dá um número finito de
mundos no modelo novo, isto é, um modelo finito. Chamemos a
parte à qual pertence um mundo~$w$ de~$[w]$. Queremos que
$\mSat{M}{!A}[w]$ se e somente se $\mSat{M^*}{!A}[{[w]}]$, pois
queremos que o novo modelo seja um contramodelo de~$!A$ se o antigo o
era. Isso exige que definamos da maneira correta tanto a partição
quanto a relação de acessibilidade de~$\mModel{M^*}$.

Para ver como isso funcionaria, imagine primeiro que não temos relação
de acessibilidade. $\mSat{M}{\Box !B}[w]$ se e somente se, para algum
$v \in W$, $\mSat{M}{\Box !B}[v]$, e o mesmo vale para
$\mModel{M^*}$, mas com $[w]$ e $[v]$. Como primeira ideia, digamos que
dois mundos $u$ e $v$ são equivalentes (pertencem à mesma parte) se
concordam quanto a todas as !!{propositional variable}s em~$\mModel{M}$,
isto é, $\mSat{M}{p}[u]$ se e somente se $\mSat{M}{p}[v]$. Seja
$V^*(p) = \Setabs{[w]}{\mSat{M}{p}[w]}$. Nosso objetivo é mostrar que
$\mSat{M}{!A}[w]$ se e somente se $\mSat{M^*}{!A}[{[w]}]$.
Evidentemente, demonstraríamos isso por indução: o caso base seria
$!A \equiv p$. Primeiro, suponha $\mSat{M}{p}[w]$. Então
$[w] \in V^*$ por definição, logo $\mSat{M^*}{p}[{[w]}]$. Suponha
agora que $\mSat{M^*}{p}[{[w]}]$. Isso significa que
$[w] \in V^*(p)$, isto é, para algum $v$ equivalente a~$w$,
$\mSat{M}{p}[v]$. Mas ``$w$ equivalente a $v$'' significa que ``$w$ e
$v$ tornam verdadeiras exatamente as mesmas !!{propositional variable}s'',
logo $\mSat{M}{p}[w]$. Para o passo indutivo, por exemplo, seja
$!A \ident \lnot !B$. Então $\mSat{M}{\lnot !B}[w]$ se e somente se
$\mSat/{M}{!B}[w]$ se e somente se $\mSat/{M^*}{!B}[{[w]}]$ (pela
hipótese de indução) se e somente se $\mSat{M^*}{\lnot !B}[{[w]}]$.
O mesmo vale para os outros operadores não modais. Isso também
funciona para $\Box$: suponha $\mSat{M^*}{\Box !B}[{[w]}]$. Isso
significa que, para todo $[u]$, $\mSat{M^*}{!B}[{[u]}]$. Pela hipótese
de indução, para todo $u$, $\mSat{M}{!B}[u]$. Consequentemente,
$\mSat{M}{\Box !B}[w]$.

No caso geral, em que também devemos definir a relação de
acessibilidade de $\mModel{M^*}$, as coisas são mais complicadas.
Chamaremos um modelo~$\mModel{M^*}$ de \emph{filtração} se sua relação
de acessibilidade~$R^*$ satisfizer as condições necessárias para que a
demonstração indutiva acima funcione. Então, qualquer
filtração~$\mModel{M^*}$ tornará $!A$ verdadeira em $[w]$ se e somente
se $\mModel{M}$ tornar $!A$ verdadeira em~$w$. Contudo, agora também
devemos mostrar que \emph{existem} filtrações, isto é, que podemos
definir~$R^*$ de modo que ela satisfaça as condições exigidas. Para que
isso funcione, contudo, devemos exigir que os mundos $u$, $v$ contem
como equivalentes não apenas quando concordam quanto a todas as
!!{propositional variable}s, mas quanto a todas as sub-!!{formula}s
de~$!A$. Como $!A$ tem apenas um número finito de sub-!!{formula}s,
isso ainda garantirá que a filtração seja finita. Não há uma única
maneira de definir uma filtração e, para garantir que a relação de
acessibilidade da filtração satisfaça as propriedades exigidas (por
exemplo, reflexividade, transitividade etc.), devemos ser criativos na
definição de~$R^*$.


\end{document}


